\xchapter{Book V, Epistle XVIII. To John, Bishop of Constantinople.}

Gregory to John, Bishop of Constantinople\footnote{On the same occasion of this letter and subsequent correspondence on the same subject, see Prolegomena, pp. xiv., xxii.}.

At the time when your Fraternity was advanced to Sacerdotal dignity, you remember what peace and concord of the churches you found. But, with what daring or with what swelling of pride I know not, you have attempted to seize upon a new name, whereby the hearts of all your brethren might have come to take offence. I wonder exceedingly at this, since I remember how thou wouldest fain have fled from the episcopal office rather than attain it. And yet, now that thou hast got it, thou desirest so to exercise it as if thou hadst run to it with ambitious intent. For, having confessed thyself unworthy to be called a bishop, thou hast at length been brought to such a pass as, despising thy brethren, to covet to be named the only bishop. And indeed with regard to this matter, weighty letters were addressed to your Holiness by my predecessor Pelagius of holy memory; in which he annulled the acts of the synod, which had been assembled among you in the case of our once brother and fellow-bishop Gregory, because of that execrable title of pride, and forbade the archdeacon whom he had sent according to custom to the threshold of our lord, to celebrate the solemnities of mass with you. But after his death, when I, unworthy, succeeded to the government of the Church, both through my other representatives and also through our common son the deacon Sabinianus, I have taken care to address your Fraternity, not indeed in writing, but by word of mouth, desiring you to restrain yourself from such presumption. And, in case of your refusing to amend, I forbade his celebrating the solemnities of mass with you; that so I might first appeal to your Holiness through a certain sense of shame, to the end that, if the execrable and profane assumption could not be corrected through shame, strict canonical measures might be then resorted to. And, since sores that are to be cut away should first be stroked with a gentle hand, I beg you, I beseech you, and with all the sweetness in my power demand of you, that your Fraternity gainsay all who flatter you and offer you this name of error, nor foolishly consent to be called by the proud title. For truly I say it weeping, and out of inmost sorrow of heart attribute it to my sins, that this my brother, who has been constituted in the grade of episcopacy for the very end of bringing back the souls of others to humility, has up to the present time been incapable of being brought back to humility; that he who teaches truth to others has not consented to teach himself, even when I implore him.

Consider, I pray thee, that in this rash presumption the peace of the whole Church is disturbed, and that it is in contradiction to the grace that is poured out on all in common; in which grace doubtless thou thyself wilt have power to grow so far as thou determinest with thyself to do so. And thou wilt become by so much the greater as thou restrainest thyself from the usurpation of a proud and foolish title: and thou wilt make advance in proportion as thou art not bent on arrogation by derogation of thy brethren. Wherefore, dearest brother, with all thy heart love humility, through which the concord of all the brethren and the unity of the holy universal Church may be preserved. Certainly the apostle Paul, when he heard some say, \emph{I am of Paul, I of Apollos, but I of Christ} (1 Cor. i. 13), regarded with the utmost horror such dilaceration of the \textsc{Lord’s} body, whereby they were joining themselves, as it were, to other heads, and exclaimed, saying, \emph{Was Paul crucified for you? or were ye baptized in the name of Paul} (ib.)? If then he shunned the subjecting of the members of Christ partially to certain heads, as if beside Christ, though this were to the apostles themselves, what wilt thou say to Christ, who is the Head of the universal Church, in the scrutiny of the last judgment, having attempted to put all his members under thyself by the appellation of Universal? Who, I ask, is proposed for imitation in this wrongful title but he who, despising the legions of angels constituted socially with himself, attempted to start up to an eminence of singularity, that he might seem to be under none and to be alone above all? Who even said, \emph{I will ascend into heaven, I will exalt my throne above the stars of heaven: I will sit upon the mount of the testament, in the sides of the North: I will ascend above the heights of the clouds; I will be like the most High} (Isai. xiv. 13).

For what are all thy brethren, the bishops of the universal Church, but stars of heaven, whose life and discourse shine together amid the sins and errors of men, as if amid the shades of night? And when thou desirest to put thyself above them by this proud title, and to tread down their name in comparison with thine, what else dost thou say but \emph{I will} \emph{ascend into heaven; I will exalt my throne above the stars of heaven?} Are not all the bishops together clouds, who both rain in the words of preaching, and glitter in the light of good works? And when your Fraternity despises them, and you would fain press them down under yourself, what else say you but what is said by the ancient foe, \emph{I will ascend above the heights of the clouds?} All these things when I behold with tears, and tremble at the hidden judgments of \textsc{God}, my fears are increased, and my heart cannot contain its groans, for that this most holy man the lord John, of so great abstinence and humility, has, through the seduction of familiar tongues, broken out into such a pitch of pride as to attempt, in his coveting of that wrongful name, to be like him who, while proudly wishing to be like \textsc{God}, lost even the grace of the likeness granted him, and because he sought false glory, thereby forfeited true blessedness. Certainly Peter, the first of the apostles, himself a member of the holy and universal Church, Paul, Andrew, John,—what were they but heads of particular communities? And yet all were members under one Head. And (to bind all together in a short girth of speech) the saints before the law, the saints under the law, the saints under grace, all these making up the \textsc{Lord’s} Body, were constituted as members of the Church, and not one of them has wished himself to be called universal. Now let your Holiness acknowledge to what extent you swell within yourself in desiring to be called by that name by which no one presumed to be called who was truly holy.

Was it not the case, as your Fraternity knows, that the prelates of this Apostolic See which by the providence of \textsc{God} I serve, had the honour offered them of being called universal by the venerable Council of Chalcedon\footnote{As to this assertion (repeated in V. 20, 43, and in VIII. 30), Giesler says, “Gregory was mistaken in believing that at the Council of Chalcedon the name \emph{Universalis Episcopus} was given to the bishop of Rome. He is styled \textgreek{οἰκουμενικὸς ἀρχιεπίσκοπος} (Mansi VI. 1006, 1012), as other patriarchs also. But in another place the title was surreptitiously introduced into the Latin acts by the Romish legates. In the sentence passed on Dioscurus, actio iii (Mansi VI. 1048), the Council say, \textgreek{ὁ ἁγιώτατος καὶ μακαριώτατος ἀρχιεπίσκοπος τῆς μεγάλης καὶ πρεσβυτέρας ῾Ρώμης Λέων}: on the contrary, in the Latin acts which Leo sent to the Gallic bishops (Leonis, Ep. 103, al. 82), we read; ‘Sanctus ac beatissimus Papa, caput universalis Ecclesiæ, Leo.’ In the older editions the beginning of Leo’s Epist. 97 (ap. Quesn. 134, Baller. 165), runs thus: ‘Leo Romæ et universalis catholicæque Ecclesiæ Episcopus Leoni semper Augusto salutem.’ Quesnel and the Ballerini, however, found in all the Codices only, ‘Leo Episcopus Leoni Augusto.’” (Giesler’s Eccl. Hist., 2nd Period, 1st Division, ch. iii. § 94, note 72).}. But yet not one of them has ever wished to be called by such a title, or seized upon this ill-advised name, lest if, in virtue of the rank of the pontificate, he took to himself the glory of singularity, he might seem to have denied it to all his brethren.

But I know that all arises from those who serve your Holiness on terms of deceitful familiarity; against whom I beseech your Fraternity to be prudently on your guard, and not to lay yourself open to be deceived by their words. For they are to be accounted the greater enemies the more they flatter you with praises. Forsake such; and, if they must needs deceive, let them at any rate deceive the hearts of worldly men, and not of priests. \emph{Let the dead bury their dead} (Luke ix. 60). But say ye with the prophet, \emph{Let them be turned back and put to shame that say unto me, Aha, Aha} (Ps. lxix. 4). And again, \emph{But let not the oil of the sinner lard my head} (Ps. cxl. 5).

Whence also the wise man admonishes well, \emph{Be in peace with many: but have but one counsellor of a thousand} (Ecclus. vi. 6). For \emph{Evil communications corrupt good manners} (1 Cor. xv. 33). For the ancient foe, when unable to break into strong hearts, looks out for weak persons who are associated with them, and, as it were, scales lofty walls by ladders set against them. So he deceived Adam through the woman who was associated with him. So, when he slew the sons of the blessed Job, he left the weak woman, that, being unable of himself to penetrate his heart, he might at any rate be able to do so through the woman’s words. Whatever weak and secular persons, then, are near you, let them be shattered in their own persuasive words and flattery, since they procure to themselves the eternal enmity of \textsc{God} from their very frowardness in being seeming lovers.

Of a truth it was proclaimed of old through the Apostle John, \emph{Little children, it is the last hour} (1 John ii. 18), according as the Truth foretold. And now pestilence and sword rage through the world, nations rise against nations, the globe of the earth is shaken, the gaping earth with its inhabitants is dissolved. For all that was foretold is come to pass. The king of pride is near, and (awful to be said!) there is an army of priests in course of preparation for him, inasmuch as they who had been appointed to be leaders in humility enlist themselves under the neck of pride. But in this matter, even though our tongue protested not at all, the power of Him who in His own person peculiarly opposes the vice of pride is lifted up for vengeance against elation. For hence it is written, \emph{\textsc{God} resisteth the proud, but giveth grace unto the humble} (Jam. iv. 6). Hence, again, it is said, \emph{Whoso exalteth his heart is unclean before \textsc{God}} (Prov. xvi. 5). Hence, against the man that is proud it is written, \emph{Why is earth and ashes proud} (Ecclus. x. 9)? Hence the Truth in person says, \emph{Whosoever} \emph{exalteth himself shall be abased} (Luke xiv. 11). And, that he might bring us back to the way of life through humility, He deigned to exhibit in Himself what He teaches us, saying, \emph{Learn of me; for I am meek and lowly in heart} (Matth. xi. 29). For to this end the only begotten Son of \textsc{God} took upon Himself the form of our weakness; to this end the Invisible appeared not only as visible but even as despised; to this end He endured the mocks of contumely, the reproaches of derision, the torments of suffering; that \textsc{God} in His humility might teach man not to be proud. How great, then, is the virtue of humility for the sake of teaching which alone He who is great beyond compare became little even unto the suffering of death! For, since the pride of the devil was the origin of our perdition, the humility of \textsc{God} has been found the means of our redemption. That is to say, our enemy, having been created among all things, desired to appear exalted above all things; but our Redeemer remaining great above all things, deigned to become little among all things.

What, then, can we bishops say for ourselves, who have received a place of honour from the humility of our Redeemer, and yet imitate the pride of the enemy himself? Lo, we know our Creator to have descended from the summit of His loftiness that He might give glory to the human race, and we, created of the lowest, glory in the lessening of our brethren. \textsc{God} humbled Himself even to our dust; and human dust sets his face as high as heaven, and with his tongue passes above the earth, and blushes not, neither is afraid to be lifted up: even man who is rottenness, and the son of man that is a worm.

Let us recall to mind, most dear brother, this which is said by the most wise Solomon. \emph{Before thunder shall go lightning, and before ruin shall the heart be exalted} (Ecclus. xxxii. 10); where, on the other hand it is subjoined, \emph{Before glory it shall be humbled.} Let us then be humbled in mind, if we are striving to attain to real loftiness. By no means let the eyes of our heart be darkened by the smoke of elation, which the more it rises the more rapidly vanishes away. Let us consider how we are admonished by the precepts of our Redeemer, who says, \emph{Blessed are the poor in spirit; for theirs is the kingdom of heaven} (Matth. v. 3). Hence, also, he says by the prophet, \emph{On whom shall my Spirit rest, but on him that is humble, and quiet, and that trembleth at my words} (Isai. lxvi. 2)? Of a truth, when the \textsc{Lord} would bring back the hearts of His disciples, still beset with infirmity, to the way of humility, He said, \emph{Whosoever will be chief among you shall be least of all} (Matth. xx. 27). Whereby it is plainly seen how he is truly exalted on high who in his thoughts is humbled. Let us, therefore, fear to be numbered among those who seek the first seats in the synagogues, and greetings in the market, and to be called of men Rabbi. For, contrariwise, the \textsc{Lord} says to His disciples, \emph{But be not ye called Rabbi: for one is your master; and all ye are brethren. And call no man your Father upon the earth, for one is your Father} (Matth. xxiii. 7, 8).

What then, dearest brother, wilt thou say in that terrible scrutiny of the coming judgment, if thou covetest to be called in the world not only father, but even general father? Let, then, the bad suggestion of evil men be guarded against; let all instigation to offence be fled from. \emph{It must needs be (indeed) that offences come; nevertheless, woe to that man by whom the offence cometh} (Matth. xviii. 7). Lo, by reason of this execrable title of pride the Church is rent asunder, the hearts of all the brethren are provoked to offence. What! Has it escaped your memory how the Truth says, \emph{Whoso shall offend one of these little ones which believe in me, it were better for him that a mill stone were hanged about his neck, and that he were drowned in the depth of the sea} (Ib. v. 6)? But it is written, \emph{Charity seeketh not her own} (1 Cor. xiii. 4). Lo, your Fraternity arrogates to itself even what is not its own. Again it is written, \emph{In honour preferring one another} (Rom. xii. 10). And thou attemptest to take the honour away from all which thou desirest unlawfully to usurp to thyself singularly. Where, dearest brother, is that which is written, \emph{Have peace with all men, and holiness, without which no man shall see the \textsc{Lord}} (Heb. xii. 14)? Where is that which is written, \emph{Blessed are the peacemakers; for they shall be called the children of \textsc{God}} (Matth. v. 9)?

It becomes you to consider, lest any root of bitterness springing up trouble you, and thereby many be defiled. But still, though we neglect to consider, supernal judgment will be on the watch against the swelling of so great elation. And we indeed, against whom such and so great a fault is committed by this nefarious attempt,—we, I say, are observing what the Truth enjoins when it says, \emph{If thy brother shall sin against thee, go and tell him his fault between thee and him alone. If he shall hear thee, thou hast gained thy brother. But if he will not hear thee, take with thee one or two more, that in the mouth of one or two witnesses every word may be established. But if he will not hear them, tell it unto the Church. But if he will not hear the Church, let him be to thee as an heathen} \emph{man and a publican} (Matth. xviii. 15). I therefore have once and again through my representatives taken care to reprove in humble words this sin against the whole Church; and now I write myself. Whatever it was my duty to do in the way of humility I have not omitted. But, if I am despised in my reproof, it remains that I must have recourse to the Church.

Wherefore may Almighty \textsc{God} show your Fraternity how great love for you constrains me when I thus speak, and how much I grieve in this case, not against you, but for you. But the case is such that in it I must prefer the precepts of the Gospel, the ordinances of the Canons, and the welfare of the brethren to the person even of him whom I greatly love.

I have received the most sweet and pleasant letter of your Holiness with respect to the case of the presbyters John and Athanasius\footnote{Cf. III. 53, and \emph{reff}.}, about which, the \textsc{Lord} helping me, I will reply to you in another letter; for, being surrounded by the swords of barbarians, I am now oppressed by such great tribulations that it is not allowed me, I will not say to treat of many things, but hardly even to breathe. Given in the Kalends of January; Indiction 13.

\xchapter{Book V, Epistle XX. To Mauricius Augustus.}

Gregory to Mauricius, \&c.

Our most pious and God-appointed lord, among his other august cares and burdens, watches also in the uprightness of spiritual zeal over the preservation of peace among the priesthood, inasmuch as he piously and truly considers that no one can govern earthly things aright unless he knows how to deal with divine things, and that the peace of the republic hangs on the peace of the universal Church. For, most serene \textsc{Lord}, what human power, and what strength of fleshly arm would presume to lift irreligious hands against the lofty height of your most Christian Empire, if the concordant hearts of priests were studious to implore their Redeemer for you with the tongue, and also, as they ought to do, by their deservings? Or what sword of a most savage race would advance with so great cruelty to the slaughter of the faithful, unless the life of us, who are called priests but are not, were weighed down by works most wicked. But while we neglect the things that concern us, and think of those that concern us not, we associate our sins with the barbaric forces and our fault, which weighs down the forces of the republic, sharpens the swords of the enemy. But what shall we say for ourselves, who press down the people of \textsc{God} which we are unworthily set over with the loads of our sins; who destroy by example what we preach with the tongue; who by our works teach unrighteous things, and with our voice only set forth the things that are righteous? Our bones are worn down by fasts, and in our mind we swell. Our body is covered with vile raiment, and in elation of heart we surpass the purple. We lie in ashes, and look down upon loftiness. Teachers of humility, we are chiefs of pride; behind the faces of sheep we hide the teeth of wolves\footnote{The ironical allusion here to John the Faster is evident.}. But what is the end of these things except that we persuade men, but are manifest to \textsc{God}? Wherefore most providently for restraining warlike movements does the most pious lord seek the peace of the Church, and, for compacting it, deigns to bring back the hearts of its priests to concord. And this indeed is what I wish; and, as far as I am concerned, I render obedience to his most serene commands. But since it is not my cause, but \textsc{God’s}, since the pious laws, since the venerable synods, since the very commands of our \textsc{Lord} Jesus Christ are disturbed by the invention of a certain proud and pompous phrase, let the most pious lord cut the place of the sore, and bind the resisting patient in the chains of august authority. For in binding up these things tightly you relieve the republic; and while you cut off such things, you provide for the lengthening of your reign.

For to all who know the Gospel it is apparent that by the \textsc{Lord’s} voice the care of the whole Church was committed to the holy Apostle and Prince of all the Apostles, Peter. For to him it is said, \emph{Peter, lovest thou Me? Feed My sheep} (John xxi. 17). To him it is said, \emph{Behold Satan hath desired to sift you as wheat; and I have prayed for thee, Peter, that thy faith fail not. And thou, when thou art converted, strengthen thy brethren} (Luke xxii. 31). To him it is said, \emph{Thou art Peter, and upon this rock I will build My Church, and the gates of hell shall not prevail against it. And I will give unto thee the keys of the kingdom of heaven and whatsoever thou shalt bind an earth shall be bound also in heaven; and whatsoever thou shalt loose on earth shall be loosed also in heaven} (Matth. xvi. 18).

Lo, he received the keys of the heavenly kingdom, and power to bind and loose is given him, the care and principality of the whole Church is committed to him, and yet he is not called the universal apostle; while the most holy man, my fellow-priest John, attempts to be called universal bishop. I am compelled to cry out and say, \emph{O tempora, O mores!}

Lo, all things in the regions of Europe are given up into the power of barbarians, cities are destroyed, camps overthrown, provinces depopulated, no cultivator inhabits the land, worshippers of idols rage and dominate daily for the slaughter of the faithful, and yet priests, who ought to lie weeping on the ground and in ashes, seek for themselves names of vanity, and glory in new and profane titles.

Do I in this matter, most pious lord, defend my own cause? Do I resent my own special wrong? Nay, the cause of Almighty \textsc{God}, the cause of the Universal Church.

Who is this that, against the evangelical ordinances, against the decrees of canons, presumes to usurp to himself a new name? Would indeed that one by himself he were, if he could be without any lessening of others,—he that covets to be universal.

And certainly we know that many priests of the Constantinopolitan Church have fallen into the whirlpool of heresy, and have become not only heretics, but even heresiarchs. For thence came Nestorius, who, thinking Jesus Christ, the Mediator of \textsc{God} and men, to be two persons, because he did not believe that \textsc{God} could be made man, broke out even into Jewish perfidy. Thence came Macedonius, who denied that \textsc{God} the Holy Spirit was consubstantial with the Father and the Son. If then any one in that Church takes to himself that name, whereby he makes himself the head of all the good, it follows that the Universal Church falls from its standing (which \textsc{God} forbid), when he who is called Universal falls. But far from Christian hearts be that name of blasphemy, in which the honour of all priests is taken away, while it is madly arrogated to himself by one.

Certainly, in honour of Peter, Prince of the apostles, it was offered by the venerable synod of Chalcedon to the Roman pontiff\footnote{Cf. V. 18, and note 5.}. But none of them has ever consented to use this name of singularity, lest, by something being given peculiarly to one, priests in general should be deprived of the honour due to them. How is it then that we do not seek the glory of this title even when offered, and another presumes to seize it for himself though not offered?

He, then, is rather to be bent by the mandate of our most pious Lords, who scorns to render obedience to canonical injunctions. He is to be coerced, who does wrong to the holy Universal Church, who swells in heart, who covets rejoicing in a name of singularity, who also puts himself above the dignity of your Empire through a title peculiar to himself.

Behold, we all suffer offence for this thing. Let then the author of the offence be brought back to a right way of life; and all quarrels of priests will cease. For I for my part am the servant of all priests, so long as they live as becomes priests. For whosoever, through the swelling of vain glory, lifts up his neck against Almighty \textsc{God} and against the statutes of the Fathers, I trust in Almighty \textsc{God} that he will not bend my neck to himself, not even with swords.

Moreover what has been done in this city on our hearing of this title, I have indicated in full to my deacon and \emph{responsalis} Sabinianus. Let then the piety of my lords think of me as their own, whom they have always cherished and countenanced beyond others, and who desire to render obedience to you and yet fear to be found guilty in the heavenly and tremendous judgment, and, according to the petition of the aforesaid deacon Sabinianus, let my most pious lord either deign to judge this business, or to move the often before mentioned man to desist at length from this attempt. If then through the most just judgment of your Piety he should comply with your orders, even though they be mild ones, we shall return thanks to Almighty \textsc{God}, and rejoice for the peace granted through you to all the Church. But should he persist any longer in his present contention, we hold this sentence of the Truth to be already made good; \emph{Every one that exalteth himself shall be humbled} (Luke xiv. 11; xviii. 14). And again it is written, \emph{Before a fall the heart is lifted up} (Prov. xvi. 18). I however, rendering obedience to the commands of my lords, have both written sweetly to my aforesaid fellow-priest, and humbly admonished him to amend himself of this coveting of empty glory. If therefore he be willing to hear me, he has a devoted brother. But, if he persists in pride, I already see what will follow:—that he will find Him as his adversary of whom it is written, \emph{\textsc{God} resisteth the proud, but giveth grace unto the humble} (Jam. iv. 6).

\xchapter{Book V, Epistle XXI. To Constantina Augusta.}

\emph{To Constantina Augusta}\footnote{The main purport of this letter to the Empress is to induce her to move the Emperor to disallow the title of Universal Bishop assumed by the patriarch of Constantinople; but at the end of the letter he takes occasion to solicit her good offices also in the case of Maximus, bishop of Salona for an account of which, with references to other letters on the subject, cf. III. 47, note 2.}.

Gregory to Constantina, \&c.

Almighty \textsc{God}, who holds in His right hand the heart of your Piety, both protects us through you and prepares for you rewards of eternal remuneration for temporal deeds. For I have learnt from the letters of the deacon Sabinianus my \emph{responsalis} with what justice your Serenity is interested in the cause of the blessed Prince of the apostles Peter against certain persons who are proudly humble and feignedly kind. And I trust in the bounty of our Redeemer that for these your good offices with the most serene lord and his most pious sons you will receive retribution also in the heavenly country. Nor is there any doubt that you will receive eternal benefits, being loosed from the chains of your sins, if in the cause of his Church you have made him your debtor to whom the power of binding and of loosing has been given. Wherefore I still beg you to allow no man’s hypocrisy to prevail against the truth, since there are some who, according to the saying of the excellent preacher, by sweet words and fair speeches seduce the hearts of the innocent,—men who are vile in raiment, but puffed up in heart. And they affect to despise all things in this world, and yet seek to acquire for themselves all the things that are of this world. They confess themselves unworthy before all men, but cannot be content with private titles, since they covet that whereby they may seem to be more worthy than all. Let therefore your Piety, whom Almighty \textsc{God} has appointed with our most serene \textsc{Lord} to be over the whole world, through your favouring of justice render service to Him from whom you have received your right to so great a dominion, that you may rule over the world that is committed to you so much the more securely as you more truly serve the Author of all things in the execution of truth.

Furthermore, I inform you that I have received a letter from the most pious lord desiring me to be pacific towards my brother and fellow-priest John. And indeed so it became the religious lord to give injunctions to priests. But, when this my brother with new presumption and pride calls himself universal bishop, having caused himself in the time of our predecessor of holy memory to be designated in synod by this so proud a title, though all the acts of that synod were abrogated, being disallowed by the Apostolic See,—the most serene lord gives me a somewhat distressing intimation, in that he has not rebuked him who is acting proudly, but endeavours to bend me from my purpose, who in this cause of defending the truth of the Gospels and Canons, of humility and rectitude; whereas my aforesaid brother and fellow-priest is acting against evangelical principles and also against the blessed Apostle Peter, and against all the churches, and against the ordinances of the Canons. But the \textsc{Lord}, in whose hands are all things, is almighty; of Him it is written, \emph{There is no wisdom nor prudence nor counsel against the \textsc{Lord}} (Prov. xxi. 30). And indeed my often before mentioned most holy brother endeavours to persuade my most serene lord of many things: but well I know that all those prayers of his and all those tears will not allow my lord to be in any thing cajoled by any one against reason or his own soul.

Still it is very distressing, and hard to be borne with patience, that my aforesaid brother and fellow-bishop, despising all others, should attempt to be called sole bishop. But in this pride of his what else is denoted than that the times of Antichrist are already near at hand? For in truth he is imitating him who, scorning social joy with the legions of angels, attempted to start up to a summit of singular eminence, saying, \emph{I will exalt my throne above the stars of heaven, I will sit upon the mount of the testament, in the sides of the North, and will ascend above the heights of the clouds, and I will be like the most High} (Isai. xiv. 13). Wherefore I beseech you by Almighty \textsc{God} not to allow the times of your Piety to be polluted by the elation of one man, nor in any way to give any assent to so perverse a title, and that in this case your Piety may by no means despise me; since, though the sins of Gregory are so great that he ought to suffer such things, yet there are no sins of the Apostle Peter that he should deserve in your times to suffer thus. Wherefore again and again I beseech you by Almighty \textsc{God} that, as the princes your ancestors have sought the favour of the holy Apostle Peter, so you also take heed both to seek it for yourselves and to keep it, and that his honour among you be in no degree lessened on account of our sins who unworthily serve him, seeing that he is able both to be your helper now in all things and hereafter to remit your sins.

Moreover, it is now even seven years that we have been living in this city among the swords of the Lombards. How much is expended on them daily by this Church, that we may be able to live among them, is not to be told. But I briefly indicate that, as in the regions of Ravenna the Piety of my Lords has for the first army of Italy a treasurer (\emph{sacellarium}) to defray the daily expenses for recurring needs, so I also in this city am their treasurer for such purposes. And yet this Church, which at one and the same time unceasingly expends so much on clergy, monasteries, the poor, the people, and in addition on the Lombards, lo it is still pressed down by the affliction of all the Churches, which groan much for this pride of one man, though they do not presume to say anything.

Further, a bishop of the city of Salona has been ordained without the knowledge of me and my \emph{responsalis}, and a thing has been done which never happened under any former princes. When I heard of this, I at once sent word to that prevaricator, who had been irregularly ordained, that he must not presume by any means to celebrate the solemnities of mass, unless we should have first ascertained from our most serene lords that they had ordered this to be done; and this I commanded him under pain of excommunication. And yet, scorning and despising me, supported by the audacity of certain secular persons, to whom he is said to give many bribes so as to impoverish his Church, he presumes up to this time to celebrate mass, and has refused to come to me according to the order of my lords. Now I, obeying the injunction of their Piety, have from my heart forgiven this same Maximus, who had been ordained without my knowledge, his presumption in passing over me and my \emph{responsalis} in his ordination, even as though he had been ordained with my authority. But his other wrong doings—to wit his bodily transgressions, which I have heard of, and his having been elected through bribery, and his having presumed to celebrate mass while excommunicated—these things, for the sake of \textsc{God}, I cannot pass over without enquiry. But I hope, and implore the \textsc{Lord}, that no fault may be found in him with respect to these things that are reported, and that his case may be terminated without peril to my soul. Nevertheless, before this has been ascertained, my most serene lord, in the order that has been despatched, has enjoined me to receive him with honour when he comes. And it is a very serious thing that a man of whom so many things of such a nature are reported should be honoured before such things have been enquired into and sifted, as they ought in the first place to be. And, if the causes of the bishops who are committed to me are settled before my most pious lords under the patronage of others, what shall I do, unhappy that I am, in this Church? But that my bishops despise me, and have recourse to secular Judges against me, I give thanks to Almighty \textsc{God} that I attribute it to my sins. This however I briefly intimate, because I am waiting for a little while; and, if he should long delay coming to me, I shall in no wise hesitate to exercise strict canonical discipline in his case. But I trust in Almighty \textsc{God}, that He will give long life to our most pious Lords, and order things for us under your hand, not according to our sins, but according to the gifts of His grace. These things, then, I suggest to my most tranquil lady, since I am not ignorant with how great zeal for rectitude the most pure conscience of her Serenity is moved.

\xchapter{Book V, Epistle XLIII. To Eulogius and Anastasius, Bishops.}

Gregory to Eulogius, Bishop of Alexandria, and Anastasius, Bishop of Antioch.

When the excellent preacher says, \emph{As long as I am the apostle of the Gentiles I will honour my ministry} (Rom. xi. 13); saying again in another place, \emph{We became as babes among you} (1 Thess. ii. 7), he undoubtedly shews an example to us who come after him, that we should retain humility in our minds, and yet keep in honour the dignity of our order, so that neither should our humility be timid nor our elevation proud. Now eight years ago, in the time of my predecessor of holy memory Pelagius, our brother and fellow-bishop John in the city of Constantinople, seeking occasion from another cause, held a synod in which he attempted to call himself Universal Bishop. Which as soon as my said predecessor knew, he despatched letters annulling by the authority of the holy apostle Peter the acts of the said synod; of which letters I have taken care to send copies to your Holiness. Moreover he forbade the deacon who attended us the most pious Lords for the business of the Church to celebrate the solemnities of mass with our aforesaid fellow-priest. I also, being of the same mind with him, have sent similar letters to our aforesaid fellow-priest, copies of which I have thought it right to send to your Blessedness, with this especial purpose, that we may first assail with moderate force the mind of our before-named brother concerning this matter, wherein by a new act of pride, all the bowels of the Universal Church are disturbed. But, if he should altogether refuse to be bent from the stiffness of his elation, then, with the succour of Almighty \textsc{God}, we may consider more particularly what ought to be done.

For, as your venerable Holiness knows, this name of Universality was offered by the holy synod of Chalcedon to the pontiff of the Apostolic See which by the providence of \textsc{God} I serve\footnote{Cf. V. 18, and note.}. But no one of my predecessors has ever consented to use this so profane a title; since, forsooth, if one Patriarch is called Universal, the name of Patriarch in the case of the rest is derogated. But far be this, far be it from the mind of a Christian, that any one should wish to seize for himself that whereby he might seem in the least degree to lessen the honour of his brethren. While, then, we are unwilling to receive this honour when offered to us, think how disgraceful it is for any one to have wished to usurp it to himself perforce.

Wherefore let not your Holiness in your epistles ever call any one Universal, lest you detract from the honour due to yourself in offering to another what is not due. Nor let any sinister suspicion make your mind uneasy with regard to our most serene lords, inasmuch as he fears Almighty \textsc{God}, and will in no way consent to do anything against the evangelical ordinances, against the most sacred canons. As for me, though separated from you by long spaces of land and sea, I am nevertheless entirely conjoined with you in heart. And I trust that it is so in all respects with your Blessedness towards me; since, when you love me in return, you are not far from me. Hence we give thanks the more to that grain of mustard seed (Matth. xiii. 31, 32), for that from what appeared a small and despicable seed it has been so spread abroad everywhere by branches rising and extending themselves from the same root that all the birds of heaven may make their nests in them. And thanks be to that leaven which, in three measures of meal, has leavened in unity the mass of the whole human race (Matth. xiii. 33); and to the little stone, which, cut out of the mountain without hands, has occupied the whole face of the earth (Dan. ii. 35), and which to this end everywhere distends itself, that from the human race reduced to unity the body of the whole Church might be perfected, and so this distinction between the several members might serve for the benefit of the compacted whole.

Hence also we are not far from you, since in Him who is everywhere we are one. Let us then give thanks to Him who, having abolished enmities, has caused that in His flesh there should be in the whole world one flock, and one sheepfold under Himself the one shepherd; and let us be ever mindful how the preacher of truth admonishes us, saying, \emph{Be careful to keep the unity of the spirit in the bond of peace} (Ephes. iv. 3), and, \emph{Follow peace with all men, and holiness, without which no man shall see \textsc{God}} (Hebr. xii. 14). And he says also to other disciples, \emph{If it be possible, as much as lieth in you, having peace with all men} (Rom. xii. 18). For he sees that the good cannot have peace with the bad; and therefore, as ye know, he premised, \emph{If it be possible.}

But, because peace cannot be established except on two sides, when the bad fly from it, the good ought to keep it in their inmost hearts. Whence also it is admirably said, \emph{As much as lieth in you;} meaning that it should remain in us even when it is repelled from the hearts of evil men. And such peace we truly keep, when we treat the faults of the proud at once with charity and with persistent justice, when we love them and hate their vices. For man is the work of \textsc{God}; but vice is the work of man. Let us then distinguish between what \textsc{God} and what man has made, and neither hate the man on account of his error nor love the error on account of the man.

Let us then with united mind attack the evil of pride in the man, that from his enemy, that is to say his error, the man himself may first be freed. Our Almighty Redeemer will supply strength to charity and justice; He will supply to us, though placed far from each other, the unity of His Spirit; even He by whose workmanship the Church, having been constructed as it were after the manner of the ark with the four sides of the world, and bound together with the compacture of incorruptible planks and the pitch of charity, is disturbed by no opposing winds, by the swelling of no billow coming from without.

But inasmuch as, with His grace steering us, we ought to seek that no wave coming upon us from without may throw us into confusion, so ought we to pray with all our hearts, dearest brethren, that the right hand of His providence may draw out the accumulation of internal bilgewater within us. For indeed our adversary the devil, who, in his rage against the humble, as a roaring lion walketh about seeking whom he may devour (1 Pet. v. 8), no longer, as we perceive, walks about the folds but so resolutely fixes his teeth in certain necessary members of the Church that, unless with the favour of the \textsc{Lord}, the heedful crowd of shepherds unanimously run to the rescue, no one can doubt that he will soon tear all the sheepfold; which \textsc{God} forbid. Consider, dearest brethren, who it is that follows close at hand, of whose approach such perverse beginnings are breaking out even in priests. For it is because he is near of whom it is written, \emph{He is king over all the sons of pride} (Job xli. 25)—not without sore grief I am compelled to say it—that our brother and fellow-bishop John, despising the \textsc{Lord’s} commands, apostolical precepts, and rules of Fathers, attempts through elation to be his forerunner in name.

But may Almighty \textsc{God} make known to your Blessedness with what sore groaning I am tormented by this consideration; that he, the once to me most modest man, he who was beloved of all, he who seemed to be occupied in alms, deeds, prayers, and fastings, out of the ashes he sat in, out of the humility he preached, has grown so boastful as to attempt to claim all to himself, and through the elation of a pompous expression to aim at subjugating to himself all the members of Christ, which cohere to one Head only, that is to Christ. Nor is it surprising that the same tempter who knows pride to be the beginning of all sin, who used it formerly before all else in the case of the first man, should now also put it before some men at the end of virtues, so as to lay it as a snare for those who to some extent seemed to be escaping his most cruel hands by the good aims of their life, at the very goal of good work, and as it were in the very conclusion of perfection.

Wherefore we ought to pray earnestly, and implore Almighty \textsc{God} with continual supplications, that He would avert this error from that man’s soul, and remove this mischief of pride and confusion from the unity and humility of the Church. And with the favour of the \textsc{Lord} we ought to concur, and make provision with all our powers, lest in the poison of one expression the living members in the body of Christ should die. For, if this expression is suffered to be allowably used, the honour of all patriarchs is denied: and while he that is called Universal perishes per chance in his error, no bishop will be found to have remained in a state of truth.

It is for you then, firmly and without prejudice, to keep the Churches as you have received them, and not to let this attempt at a diabolical usurpation have any countenance from you. Stand firm; stand secure; presume not ever to issue or to receive writings with the falsity of the name Universal in them. Bid all the bishops subject to your care abstain from the defilement of this elation, that the Universal Church may acknowledge you as Patriarchs not only in good works but also in the authority of truth. But, if perchance adversity is the consequence, we ought to persist unanimously, and show even by dying that in case of harm to the generality we do not love anything of our own especially. Let us say with Paul, \emph{To me to live is Christ, and to die is gain} (Philip. i. 21). Let us hear what the first of all pastors says; \emph{If ye suffer anything for righteousness’ sake, happy are ye} (1 Pet. iii. 14). For believe me that the dignity which we have received for the preaching of the truth we shall more safely relinquish than retain in behalf of the same truth, should case of necessity require it. Finally, pray for me, as becomes your most dear Blessedness, that I may shew forth in works what I am thus bold to say to you.

\xchapter{Book VII, Epistle XXVII. To Anastasius, Bishop.}

Gregory to Anastasius, Bishop of Antioch.

I have received through the hands of our common son the deacon Sabinianus the longed for letter of your most sweet Holiness, in which the words have flowed not from your tongue but from your soul. And it is not surprising that one speaks well who lives perfectly. And, since you have learnt, through the Spirit teaching you in the school of the heart, the precepts of life—to despise all earthly things and to speed to the heavenly country,—in proportion as you have advanced in good you think what is good of others. But, when I heard many things said in the letters of your Blessedness in praise of me, I understood your intention; how that you wished to describe not what I am, but what I ought to be. But as to your saying that I ought to remember my manner of life, and on no account give place to the malignant spirit who seeks to sift souls, I indeed recollect myself to have been always of bad manner of life, and hasten to overcome and put an end to this my manner of life, if I can. If however, as you believe, I have had anything good in me, I trust in the help of Almighty \textsc{God} that I have not forgotten it. But your Holiness, as I see, by the words of sweetness at the beginning and the words that follow, has wished your letter to be like a bee, which carries both honey and a sting, satiating me with the honey and piercing me with the sting. But meanwhile I return to meditation on the words of Solomon, \emph{That better are the wounds of one that loves than the kisses of a flattering foe} (Prov. xxvii. 6). Thus, as to your saying that we ought not to give occasion of offence for no cause at all, this is what your son, our most pious lord (for whose life we ought continually to pray) has already written repeatedly; and what he says out of power I know that you say out of love. Nor do I wonder that you have made use of imperial language in your letters, since there is a very close relationship between love and power. For both presume in a princely way; both ever speak with authority.

And indeed on the receipt of the synodical epistle of our brother and fellow-bishop Cyriacus it was not worth my while to make a difficulty on account of the profane title at the risk of disturbing the unity of holy Church: but nevertheless I took care to admonish him with respect to this same superstitious and proud title, saying that he could not have peace with us unless he corrected the elation of the aforesaid expression, which the first apostate invented. You, however, ought not to say that this is a matter of no consequence, since, if we bear it with equanimity, we are corrupting the faith of the Universal Church; for you know how many not only heretics but heresiarchs have issued from the Constantinopolitan Church. And, not to speak of the injury done to your dignity, if one bishop is called Universal, the Universal Church comes to ruin, if the one who is universal falls. But far, far be this levity from my ears. Yet I trust in Almighty \textsc{God} that what He has promised He will soon fulfil; \emph{Whosoever exalteth himself shall be humbled} (Luke xiv. 11).

So much, in the midst of many occupations. I have briefly replied to what you have said in your letters: for what I ought not just now to express in writing remains imprinted on my mind. I beg your Blessedness always to recall me to your memory in your holy prayers, that so your intercessions may rescue me from temporal and eternal ills. Pray moreover zealously and fervently for the most serene lord the Emperor; for his life is very necessary for the world. I refrain from saying more, for I doubt not that you know.

\xchapter{Book VIII, Epistle XXX. To Eulogius, Bishop of Alexandria.}

To Eulogius, Bishop of Alexandria.

Gregory to Eulogius, \&c.

Our common son, the bearer of these presents, when he brought the letters of your Holiness found me sick, and has left me sick; whence it has ensued that the scanty water of my brief epistle has been hardly able to exude to the large fountain of your Blessedness. But it was a heavenly boon that, while in a state of bodily pain, I received the letter of your Holiness to lift me up with joy for the instruction of the heretics of the city of Alexandria, and the concord of the faithful, to such an extent that the very joy of my mind moderated the severity of my suffering. And indeed we rejoice with new exultation to hear of your good doings, though at the same time we by no means suppose that it is a new thing for you to act thus perfectly. For that the people of holy Church increases, that spiritual crops of corn for the heavenly garner are multiplied, we never doubted that this was from the grace of Almighty \textsc{God} which flowed largely to you, most blessed ones. We therefore render thanks to Almighty \textsc{God}, that we see fulfilled in you what is written, \emph{Where there is much increase, there the strength of the oxen is manifest} (Prov. xiv. 4). For, if a strong ox had not drawn the plough of the tongue over the ground of the hearts of hearers, so great an increase of the faithful would by no means have sprung up.

But, since in the good things you do I know that you also rejoice with others, I make you a return for your favour, and announce things not unlike yours; for while the nation of the Angli, placed in a corner of the world, remained up to this time misbelieving in the worship of stocks and stones, I determined, through the aid of your prayers for me, to send to it, \textsc{God} granting it, a monk of my monastery for the purpose of preaching. And he, having with my leave been made bishop by the bishops of Germany, proceeded, with their aid also, to the end of the world to the aforesaid nation; and already letters have reached us telling us of his safety and his work; to the effect that he and those that have been sent with him are resplendent with such great miracles in the said nation that they seem to imitate the powers of the apostles in the signs which they display. Moreover, at the solemnity of the \textsc{Lord’s} Nativity which occurred in this first indiction, more than ten thousand Angli are reported to have been baptized by the same our brother and fellow-bishop. This have I told you, that you may know what you are effecting among the people of Alexandria by speaking, and what in the ends of the world by praying. For your prayers are in the place where you are not, while your holy operations are shewn in the place where you are.

In the next place, as to the person of Eudoxius the heretic\footnote{Cf. VII. 4, and 34.}, about whose error I have discovered nothing in the Latin language, I rejoice that I have been most abundantly satisfied by your Blessedness. For you have adduced the testimonies of the strong men, Basil, Gregory, and Epiphanius; and we acknowledge him to be manifestly slain, at whom our heroes have cast so many darts. But with regard to these errors which are proved to have arisen in the Church of Constantinople, you have replied on all heads most learnedly, and as it became you to utter the judgment of so great a see. Whence we give thanks to Almighty \textsc{God}, that the tables of the covenant are still in the ark of \textsc{God}. For what is the priestly heart but the ark of the covenant? And since spiritual doctrine retains its vigour therein, without doubt the tables of the law are lying in it.

Your Blessedness has also been careful to declare that you do not now make use of proud titles, which have sprung from a root of vanity, in writing to certain persons, and you address me saying, \emph{As you have commanded.} This word, \emph{command}, I beg you to remove from my hearing, since I know who I am, and who you are. For in position you are my brethren, in character my fathers. I did not, then, command, but was desirous of indicating what seemed to be profitable. Yet I do not find that your Blessedness has been willing to remember perfectly this very thing that I brought to your recollection. For I said that neither to me nor to any one else ought you to write anything of the kind; and lo, in the preface of the epistle which you have addressed to myself who forbade it, you have thought fit to make use of a proud appellation, calling me Universal Pope. But I beg your most sweet Holiness to do this no more, since what is given to another beyond what reason demands is subtracted from yourself. For as for me, I do not seek to be prospered by words but by my conduct. Nor do I regard that as an honour whereby I know that my brethren lose their honour. For my honour is the honour of the universal Church: my honour is the solid vigour of my brethren. Then am I truly honoured when the honour due to all and each is not denied them. For if your Holiness calls me Universal Pope, you deny that you are yourself what you call me universally. But far be this from us. Away with words that inflate vanity and wound charity.

And, indeed, in the synod of Chalcedon and afterwards by subsequent Fathers, your Holiness knows that this was offered to my predecessors\footnote{Cf. V. 18, note 5.}. And yet not one of them would ever use this title, that, while regarding the honour of all priests in this world, they might keep their own before Almighty \textsc{God}. Lastly, while addressing to you the greeting which is due, I beg you to deign to remember me in your holy prayers, to the end that the \textsc{Lord} for your intercessions may absolve me from the bands of my sins, since my own merits may not avail me.
